(1) 旧基态满足高斯波函数 \(\psi_i(x)=(m\omega_i/(\pi\hbar))^{1/4}\exp[-m\omega_ix^2/(2\hbar)]\)。标准高斯积分给位置方差 \(\langle x^2\rangle_i=\hbar/(2m\omega_i)\)，动量方差 \(\langle p^2\rangle_i=m\hbar\omega_i/2\)，均值 \(\langle x\rangle=\langle p\rangle=0\)。瞬间改变势参数时状态连续，因此这些期望值在 \(t=0^+\) 仍与 \(t=0^-\) 相同。
(2) 新哈密顿量 \(H_f=p^2/(2m)+m\omega_f^2x^2/2\)，所以 \(\langle H_f\rangle=\hbar(\omega_i+\omega_f^2/\omega_i)/4\)。由 \(\langle H_f\rangle=\hbar\omega_f(\langle n_f\rangle+1/2)\) 得 \(\langle n_f\rangle=\frac14(\omega_i/\omega_f+\omega_f/\omega_i-2)=(\omega_f-\omega_i)^2/(4\omega_i\omega_f)\ge0\)。不等式保证突然过程不会把旧基态变成能量低于新基态的状态。
(3) 两个正归一化高斯基态的振幅重叠由积分 \(\int_{-\infty}^\infty\psi_f^*(x)\psi_i(x)\,dx\) 给出，模平方为 \(P_0=2\sqrt{\omega_i\omega_f}/(\omega_i+\omega_f)\)，在频率不变时为一，其余情况小于一。旧基态波函数关于 \(x\) 为偶函数，新谐振子的奇数本征态为奇函数，因此对应重叠积分严格为零；骤变只激发偶数能级。
(4) 瞬间外部参数改变的平均功为 \(W=\langle H_f\rangle-\hbar\omega_i/2=\hbar(\omega_f^2-\omega_i^2)/(4\omega_i)\)，可正可负，负值表示降低势刚度时外部提取能量。若频率变化足够缓慢且系统保持隔离并满足量子绝热条件，起始基态会演化为新基态，最终能量仅为 \(\hbar\omega_f/2\)；突然淬火额外增加的激发能是 \(\hbar(\omega_f-\omega_i)^2/(4\omega_i)\)。
